\section*{Appendix}

In addition to the code blocks attached in the Appendix, the complete code can be found at the
following URL: \href{https://github.com/vbhaga89/MIT801_ML_Exam}{GitHub Repository}

\subsection*{Question 3.2 c}

\begin{lstlisting}[language=Python, caption={Grid search for RBF-SVM hyperparameter tuning}]
from sklearn.model_selection import GridSearchCV, StratifiedKFold

# Define values to search C and gamma
param_grid = {
    'C': [0.1, 1, 10],
    'gamma': [0.01, 0.1, 1]
}

# Create model
svm = SVC(
    kernel="rbf",
    probability=True,
    random_state=42
)

# Define cross validation with stratification (ensures class balance in splits)
cv = StratifiedKFold(
    n_splits=5,
    shuffle=True,
    random_state=42
)

# Create grid search class using AUC of ROC curve as scoring metric
grid = GridSearchCV(
    estimator=svm,
    param_grid=param_grid,
    scoring='roc_auc',
    cv=cv,
    n_jobs=-1
)

grid.fit(X_train, y_train)
print("Best parameters:", grid.best_params_)
\end{lstlisting}

\newpage

\subsection*{Question 4.2}

\begin{lstlisting}[language=Python, caption={MSE, cross-entropy and hinge loss functions' calculations and plots}]
import numpy as np
import matplotlib.pyplot as plt

# Prediction scores from -2 to 2
y_pred = np.linspace(-2, 2, 500)

# Positive example
y_true = 1

# Sigmoid function for cross-entropy
def sigmoid_func(z):
    return 1 / (1 + np.exp(-z))

# Loss functions
mse = (y_pred - y) ** 2
cross_entropy = -np.log(sigmoid_func(y_pred))
hinge = np.maximum(0, 1 - y * y_pred)

# Plot
plt.figure(figsize=(8, 5))

plt.plot(y_pred, mse, label="MSE Loss")
plt.plot(y_pred, cross_entropy, label="Cross-Entropy Loss")
plt.plot(y_pred, hinge, label="Hinge Loss")

plt.xlabel(r"Prediction score ($\hat{y})$")
plt.ylabel("Loss")
plt.title("MSE, Cross-Entropy, and Hinge Loss for a Positive Example (y = 1)")
plt.legend()
plt.grid(True)

plt.show()
\end{lstlisting}

\newpage

\subsection*{Question 5.1 b.}

\begin{lstlisting}[language=Python, caption={PCA and Explained Variance Ratio plot}]
from sklearn.decomposition import PCA

# Fit PCA on scaled training data
pca = PCA()
X_train_pca = pca.fit_transform(X_train)

# Explained variance ratios
explained_var = pca.explained_variance_ratio_

# Get cumulative variances
cumulative_var = np.cumsum(explained_var)

# Minimum number of components needed for at least 90% variance
n_components_90 = np.argmax(cumulative_var >= 0.90) + 1

print("Explained variance ratio:", explained_var)
print("Cumulative explained variance:", cumulative_var)
print("Number of components for at least 90% variance:", n_components_90)

# Plot explained variance ratio
plt.figure(figsize=(8, 5))

# Plot the explained variance of each component
plt.bar(
    range(1, len(explained_var) + 1),
    explained_var,
    label="Individual explained variance",
    color='orange'
)

# Plot the cumalative variance as you add components
plt.plot(
    range(1, len(cumulative_var) + 1),
    cumulative_var,
    marker="o",
    label="Cumulative explained variance"
)

# Plot threshold line
plt.axhline(y=0.90, linestyle="--", color='red', label="90% threshold")

plt.xlabel("Principal component")
plt.ylabel("Explained variance ratio")
plt.title("PCA Explained Variance Ratio")
plt.xticks(range(1, len(explained_var) + 1))
plt.legend()
plt.grid(True)
plt.show()
\end{lstlisting}

\newpage

\subsection*{Question 6.2}

\begin{lstlisting}[language=Python, caption={ROC Curve and AUC for all 3 models}]
from sklearn.metrics import roc_curve, roc_auc_score
from sklearn.linear_model import LogisticRegression

# Define Models
sgd_logistic = SGDClassifier(
    loss="log_loss",
    learning_rate="optimal",
    max_iter=200,
    random_state=42
)
rbf_svm = SVC(
    kernel="rbf",
    C=1.0,
    gamma="scale",
    probability=True,
    random_state=42
)
ridge_logistic = LogisticRegression(
    penalty='l2',
    solver='lbfgs',
    max_iter=1000,
    random_state=42
)

models = {
    "SGD Logistic Regression": sgd_logistic,
    "RBF-SVM": rbf_svm,
    "Ridge Logistic Regression": ridge_logistic
}

# Fit the Models and Plot ROC with AUC score
plt.figure(figsize=(8, 6))

for name, model in models.items():
    model.fit(X_train, y_train)
    # Get predicted probalities
    y_test_proba = model.predict_proba(X_test)[:, 1]

    # Calculate ROC curve and AUC
    fpr, tpr, thresholds = roc_curve(y_test, y_test_proba) 
    roc_auc = roc_auc_score(y_test, y_test_proba)

    # Plot ROC curve
    plt.plot(fpr, tpr, label=f"{name} AUC = {roc_auc:.3f}")

# Add Baseline of AUC = 0.5 curve
plt.plot([0, 1], [0, 1], linestyle="--", label="Random Classifier AUC = 0.500")
# Plot the ROC curve
plt.xlabel("False Positive Rate")
plt.ylabel("True Positive Rate")
plt.title("ROC Curve Comparison for Churn Models")
plt.legend()
plt.grid(True)
plt.show()
\end{lstlisting}